\section{Discussion and Limitations}\label{sec:discussion}
The three questions distinguish a capability from its coverage and mechanism.
RQ1 establishes shared tracking of held-out motions on selected trained
bodies and execution across all trained bodies on development clips. RQ2
establishes high average transfer on both sides of the training coefficient
boundary, while exposing bodies for which tracking largely fails. RQ3
supports practical architecture and reference-processing choices, but not a
causal explanation of transfer.

\paragraph{Average transfer and individual robustness.}
Similar inside/outside averages do not make body extrapolation uniformly
reliable. Tier success is non-monotonic, and poor tracking occurs both beyond
and within the training mass/height ranges. These simulated properties
describe the tested assets; their associations do not identify failure causes.
The later checkpoint recovers some bodies and substantially harms others.
Its outside success gain accompanies a slight increase in mean position
error, so threshold success alone would overstate the improvement.

\paragraph{Coverage and inference.}
The unseen-body headline uses known motion identities to isolate body
transfer. The separate 23 validation/test clips cannot provide an untouched
motion test for the full-data checkpoint, which trained on them. The
1,048-clip independent test uses one trained body per clip; a fixed
held-out-motion panel crossed with all 128 trained bodies remains missing.
Original/mirrored variants stay together, but the splits are not demonstrated
to be subject- or recording-disjoint. Clip-bootstrap intervals hold the body
panel fixed and do not quantify uncertainty over shape space. One full-scale
training run and one candidate checkpoint do not measure training-seed or
ordinary checkpoint variation. Identical-checkpoint reruns measure simulator
variability, with nearly cancelling pooled changes but appreciable pair flips.

\paragraph{Mechanism and design.}
State and reference observations already carry shape information. A matched
policy without explicit morphology inputs is needed to attribute transfer to
those inputs. The architecture comparison keeps them in every model and has
unequal capacities. Refinement reduces contact artifacts with small pose
changes but raises target jerk; policies evaluated against different targets
do not isolate its learning benefit. The exploratory physical analysis in
Appendix~\ref{sec:physics} remains descriptive pending a clip-level contrast.

\paragraph{Failure diagnosis and scope.}
The unseen-body runs detect error-threshold failures, and all failed rows have zero termination flags. Those flags alone do not
prove absence of falls; their interpretation depends on the replay
configuration. Synchronized rollouts, contact/root trajectories, and reference-quality
audits are needed before assigning a failure mechanism. Extreme HUMOS
references have been refined but lack per-asset feasibility validation.
Estimated masses, collision proxies, and mass-scaled gains are modeling
choices; tracking success and jerk do not certify biomechanical realism.
The short, reference-initialized flat-ground protocol excludes objects,
partners, long-horizon recovery, skeletal-topology changes, and real robots.

The corpus is derived: clip/body variants are not independent captures.
Caption alignment, immutable generation checksums, and redistribution
permissions require auditing before a benchmark-release claim. A matched
external tracker would strengthen comparative conclusions, but the published
capability table alone cannot establish a performance advantage.
